% !TeX root = ../深度学习笔记.tex
% 由主文档通过 \input 载入；请编译主文档。
\setcounter{chapter}{16} % 保留原书章号。
\chapter{语言建模与生成}
\label{chap:dl-language}
语言模型把文本映成离散序列并预测条件分布。训练学习的是概率规律，生成还需要选择规则；高概率文本、符合指令的回答和事实正确的回答并非同一个概念。

\section{分词、词表与序列数据}
\emph{分词（tokenization）}把文本变成词元编号序列。字符词表小但序列较长，整词词表较大且存在未见词，子词方法在二者之间组合常见片段。中文的字符、词语与子词切分也可能不同；同一文本在不同分词器下的长度不能直接比较。特殊词元用于标记开始、结束、填充或角色，它们是模型输入协议的一部分。

训练窗口可以由长文档切片得到，但必须保留需要的文档边界规则。若跨文档连接，模型实际学习了何种连接分布需要解释；若相邻窗口重叠，评价划分应避免同一来源片段泄漏。词表与规范化步骤应在预定训练信息范围内确定，评价时使用相同编码和解码规则。

\section{自回归目标与困惑度}
对固定长度的词元序列 $x_{1:T}$，条件概率链式法则给出
\begin{equation}
 p_\theta(x_{1:T})=\prod_{t=1}^Tp_\theta(x_t\mid x_{<t}),\qquad
 \operatorname{PPL}=\exp\!\left(-\frac1T\sum_{t=1}^T\log p_\theta(x_t\mid x_{<t})\right).
\label{eq:dl-language-probability}
\end{equation}
其中首位置使用起始上下文，\emph{困惑度（perplexity）}是平均负对数似然的指数。对整个语料应先合并有效词元的负对数似然，再取指数，而不是无权平均各句困惑度。可变长度完整文本还须对结束词元建模，或明确给定长度；截断生成不等于模型赋予截断文本全部的序列概率。

训练通常对真实前缀执行教师强制，最小化有效词元交叉熵。若模型始终给正确词元概率 $1/K$，困惑度为 $K$，可理解为这类均匀预测器的等效候选数。但不同分词单位、语料和边界协议改变这一量的尺度，跨设置的困惑度不可直接作优劣结论。

\section{解码策略与生成行为}
贪心解码每步取最大概率词元；束搜索保留若干高分前缀并扩展，是完整序列搜索的近似。温度采样使用 $p_k^{(\tau)}=e^{z_k/\tau}/\sum_je^{z_j/\tau}$，$\tau>0$，温度改变分布尖锐程度；top-$k$ 在最高分的 $k$ 个候选内重归一化，top-$p$ 则保留累计质量达到指定阈值的高概率集合。生成分布因这些选择而可能不同于原模型分布。

贪心不保证全序列最优。例如首步 $P(A)=0.6,P(B)=0.4$，在 $A$ 后最佳下一词概率为 $0.51$，在 $B$ 后为 $0.99$；贪心得到两步概率 $0.306$，而从 $B$ 开始可得 $0.396$。另一方面，序列概率最优也不等于任务答案最好。生成评价还需检查正确性、约束满足、重复与结束行为，采样时记录随机种子和解码参数。

\section{指令微调、上下文学习与检索增强}
\emph{指令微调（instruction tuning）}用指令、可选上下文和回答构造监督样本。对条件文本 $c$ 与回答 $a_{1:T}$，常见目标是 $-\sum_t\log p_\theta(a_t\mid c,a_{<t})$，损失只施加在需要学习生成的回答位置；输入中保留的提示不等于都要计入目标。偏好数据则提供回答间的比较，如何转成目标还需另外规定奖励或偏好优化模型。

\emph{上下文学习（in-context learning）}把示例或规则写入输入，在固定参数下改变条件分布；它不等于通过梯度微调。\emph{检索增强生成（retrieval-augmented generation, RAG）}先取回相关材料，再把材料作为条件输入。检索缺失、证据错误和生成不忠于证据会产生不同失败，须分别评价。文本流畅或模型自述“有把握”均不能替代事实核查与任务标准。

\section{例题与习题}
\begin{example}[计算词元困惑度]
两个目标词元在真实前缀下的条件概率分别为 $1/2$ 和 $1/4$。求序列概率、平均负对数似然与困惑度。
\end{example}
\begin{solve}
序列概率为 $1/8$，平均负对数似然为 $(\log2+\log4)/2=\frac32\log2$，困惑度为 $2^{3/2}=\sqrt8\approx2.828$。两个位置概率的算术平均为 $3/8$，取其倒数不等于正确困惑度，因为该指标来自对数平均。
\end{solve}
\begin{enumerate}
\item 在本章两步概率树上比较贪心与宽度二的束搜索，补全每个结点剩余候选概率。
\item 训练小型字符模型，分别报告困惑度与固定提示下的生成样本，检查是否记忆训练片段。
\item 设计固定参数、有无检索材料的对照，并把检索命中和回答正确分别计分，说明能据此识别哪些错误。
\end{enumerate}

